\section[Method]{The method}
\sectionframe{02}{The method}{How does the method keep edges sharp?}{\faCogs}

% ------------------------------------------------------------------ pipeline
% Blocks from left to right with what flows along each arrow; the block under discussion in colour.
\begin{frame}[source={Author et al.\ (2026), Journal 1, 1, after Fig.\ 1}]{The pipeline}
  \begin{tikzpicture}[x=1mm,y=1mm]
    \useasboundingbox (0,0) rectangle (144,-54);
    % input: panel a of the example figure; output: panel c
    \figpart{0}{-6}{assets/example-figure}{4000}{0}{0}{1000}{1000}{28}
    \node[small label,anchor=north] at (14,-35.2) {input};
    \draw[flow] (29.4,-20) -- node[lbl,above] {image} (39,-20);
    \node[box,minimum width=20mm,minimum height=11mm] at (49.6,-20) {encoder};
    \draw[flow] (60.2,-20) -- node[lbl,above] {512 numbers} (76.4,-20);
    \node[primary box,minimum width=20mm,minimum height=11mm] at (87,-20) {decoder};
    \draw[flow] (97.6,-20) -- node[lbl,above] {image} (106.6,-20);
    \figpart{108}{-6}{assets/example-figure}{4000}{2000}{0}{3000}{1000}{28}
    \node[small label,anchor=north] at (122,-35.2) {output};
    % training: the loss compares the output with the input and updates the decoder
    \draw[accent,line width=0.6pt,dash pattern=on 1.6pt off 1.2pt] (14,-38.6) -- (14,-45) -- (122,-45)
      -- (122,-38.6);
    \draw[accent,line width=0.6pt,dash pattern=on 1.6pt off 1.2pt,-{Stealth[length=1.6mm,width=1.2mm]}]
      (87,-45) -- (87,-26.2);
    \node[anchor=north,small label,text=accent] at (49,-46.2) {training: the loss compares output and input};
  \end{tikzpicture}
  \note{\textbf{2:10--3:10}\par
    The model is an encoder and a decoder: the encoder turns an image into 512 numbers, the decoder turns
      them back into an image.\par
    Training compares each output with its input, the red dashed path, and updates the model.\par
    The method changes only the loss, the next slide; the decoder is where its effect shows.}
\end{frame}

% ---------------------------------------------------------------------- loss
% The equation large and alone, each term coloured by its meaning and labelled.
\begin{frame}[source={Author et al.\ (2026), Journal 1, 1, Eqs.\ 1--3 · Other and Author (2020), Conference}]
  {The loss}
  \begin{tikzpicture}[x=1mm,y=1mm]
    \useasboundingbox (0,0) rectangle (144,-54);
    \node[font=\LARGE,text=ink,inner sep=0pt] at (36,-8) {$L\;=$};
    \node[font=\LARGE,text=grey,inner sep=0pt] (data) at (55,-8) {$L_{\mathrm{data}}$};
    \node[font=\LARGE,text=ink,inner sep=0pt] at (70,-8) {$+$};
    \node[font=\LARGE,text=accent,inner sep=0pt] (lam) at (80,-8) {$\lambda$};
    \node[font=\LARGE,text=primary,inner sep=0pt] (edge) at (95,-8) {$L_{\mathrm{edge}}$};
    % one column per term: what it compares, then its formula
    \foreach \x/\c/\t in {40/grey/{per pixel: output against input; the baseline's whole loss},
        82/accent/{weight of the edge term},
        124/primary/{image gradients: large along edges}}
      \node[anchor=north,label text,text=\c,align=center,text width=36mm] at (\x,-24) {\t};
    \draw[flow,draw=grey] (40,-23) -- (data.south);
    \draw[flow,draw=accent] (82,-23) -- (lam.south);
    \draw[flow,draw=primary] (124,-23) -- (edge.south);
    \node[anchor=north,font=\small,text=ink,inner sep=0pt] at (40,-37)
      {$\displaystyle\frac{1}{n}\sum_{i=1}^{n}|x_i-\hat{x}_i|$};
    \node[anchor=north,font=\small,text=ink,inner sep=0pt] at (82,-39.4) {$\lambda=0.1$};
    \node[anchor=north,font=\small,text=ink,inner sep=0pt] at (124,-37)
      {$\displaystyle\frac{1}{n}\sum_{i=1}^{n}|\nabla x_i-\nabla\hat{x}_i|$};
  \end{tikzpicture}
  \note{\textbf{3:10--4:10}\par
    The loss adds one term to the baseline's.\par
    The data term compares output and input pixel by pixel; alone, it is the baseline's loss.\par
    The edge term compares their gradients, which are large only along edges, so a blurred edge now
      costs something.\par
    The weight lambda sets how much; the next slide shows what it changes.}
\end{frame}

% ------------------------------------------------------------------- weights
% Options side by side as cards with the same rows: a head, the same image region, its properties.
\begin{frame}[source={Author et al.\ (2026), Journal 1, 1, Fig.\ 3}]{Choosing the weight}
  \begin{tikzpicture}[x=1mm,y=1mm]
    \useasboundingbox (0,0) rectangle (144,-54);
    % the same region of panels b, c and d: the two squares and their shared edge
    \foreach \k/\c/\h/\p/\one/\two in {
        0/greymid/{$\lambda=0$, the baseline}/1000/{edges blur}/{lowest pixel error},
        1/primary/{$\lambda=0.1$}/2000/{edges sharp}/{pixel error close to the baseline's},
        2/accent/{$\lambda=1$}/3000/{halos along edges}/{grain in flat areas}} {
      \cardhead{\k*49}{0}{46}{\c}{\h}
      \fill[inktint] (\k*49,-5.4) rectangle ++(46,-48.6);
      \figpart{\k*49+2}{-7.4}{assets/example-figure}{4000}{\p+540}{430}{\p+980}{710}{42}
      \node[card text,text width=42mm] at (\k*49+2,-36.6) {\cardbullet{\one}\cardbullet{\two}};}
  \end{tikzpicture}
  \note{\textbf{4:10--5:10}\par
    The same region, the two squares, for three weights.\par
    Without the edge term the edge blurs; with a weight of 0.1 it stays sharp at almost the same pixel
      error; with a weight of 1, halos and grain appear.\par
    The weight was chosen on the validation set, not on the test set.}
\end{frame}

% ----------------------------------------------------------------- running it
% Commands in a terminal window, the files they use as a folder tree, and one file shown as in an editor.
\begin{frame}[fragile,source={Author et al.\ (2026), code repository}]{Running it}
  \begin{columns}[T,onlytextwidth]
    \begin{column}{0.56\textwidth}
\begin{terminal}{training and evaluation}
$ python train.py --config configs/edge.yaml
epoch  1/50   loss 0.412
epoch 50/50   loss 0.087
$ python evaluate.py --split test
error 0.27   baseline 0.42
\end{terminal}
      \vspace{2mm}
      \begin{foldertree}[note at=30,row height=4]
        \tdir{0}{project/}{}
        \tdir{1}{configs/}{}
        \tfile{2}{edge.yaml}{the loss and its weight}
        \tfile{1}{train.py}{trains the model}
        \tfile{1}{evaluate.py}{the error on a split}
      \end{foldertree}
    \end{column}
    \begin{column}{0.4\textwidth}
\begin{codefile}[style=sh]{configs/edge.yaml}
# the loss of slide 4
loss:
  data: l1
  edge: sobel
  weight: 0.1
\end{codefile}
    \end{column}
  \end{columns}
  \note{\textbf{5:10--6:00}\par
    Training takes one command and one configuration file.\par
    The configuration names the two loss terms and the weight; everything else is the baseline's.\par
    Evaluation runs on the test set, which training never sees.}
\end{frame}
